\documentclass[11pt]{article}

\usepackage[margin=1in]{geometry}
\usepackage{amsmath}
\usepackage{amssymb}
\usepackage{booktabs}
\usepackage[destlabel]{hyperref}

\hypersetup{colorlinks=true, linkcolor=blue, citecolor=blue}

\title{Fund-Flow Rotation Across Equity Sectors from Public Filings}
\author{Max Piterbarg}
\date{\today}

\begin{document}
\maketitle

\begin{abstract}
We measure where investor capital rotates across the eleven sectors of the US equity
market, separating investor demand from price performance, using only public data. The
input is the net dollar creation and redemption of exchange-traded fund shares, which
each fund reports monthly on SEC Form N-PORT. We normalize these flows by fund size,
aggregate to sectors, and compare each sector both to the market as a whole and to its
own history. The result is a monthly sector-flow panel covering 2019 to 2026 whose
reported returns match price-derived returns at a correlation above $0.99$, and whose
flows and returns reconcile with reported net assets to within about one percent per
quarter. On that panel we define the standard rotation statistics: cumulative flow,
standardized flow, relative flow, and flow momentum.
\end{abstract}

\section{Flow versus price}
\label{sec:flow-vs-price}

A fund's assets under management move for two reasons. The market reprices the holdings
the fund already owns, and investors add or withdraw money. Only the second is a flow.
Over one period the two combine as
\begin{equation}
\label{eq:aum_identity}
A_t = A_{t-1}\,(1 + r_t) + F_t,
\end{equation}
where $A_t$ is assets under management at the end of period $t$, $r_t$ is the fund's
total return over the period, and $F_t$ is the net dollar flow. The first term is what
the previous period's money grew to on its own. Anything beyond it is new investor money.
Rearranging \eqref{eq:aum_identity} isolates the quantity of interest,
\begin{equation}
\label{eq:flow_residual}
F_t = A_t - A_{t-1}\,(1 + r_t).
\end{equation}

A flow is a demand signal, not a price signal. Consider a fund worth $100$ at the start
of a day. The market rises two percent, lifting its holdings to $102$ with no investor
action. At the close the fund is worth $101$. Assets rose, from $100$ to $101$, yet
\eqref{eq:flow_residual} gives $F = 101 - 102 = -1$: investors withdrew one unit while
the price rose. Assets under management and flow can move in opposite directions, which
is what makes flow an independent measure of what investors chose to do.

\section{Reading the flow}
\label{sec:reading-flow}

Equation \eqref{eq:flow_residual} requires assets and return, which still leaves $F_t$
inferred rather than observed. Two routes make it more direct.

\subsection{The shares route}

An exchange-traded fund issues new shares when demand arrives and redeems them when
demand leaves, each at approximately net asset value. The change in shares outstanding
therefore reveals the flow,
\begin{equation}
\label{eq:flow_shares}
F_t = (S_t - S_{t-1})\,\mathrm{NAV}_t,
\end{equation}
with $S_t$ shares outstanding and $\mathrm{NAV}_t$ net asset value per share at the end
of $t$. This is \eqref{eq:flow_residual} specialized to an ETF: writing $A_t = S_t\,
\mathrm{NAV}_t$ and $r_t = \mathrm{NAV}_t / \mathrm{NAV}_{t-1} - 1$ and substituting
recovers \eqref{eq:flow_shares} exactly.

A share split breaks \eqref{eq:flow_shares}. A two-for-one split doubles $S_t$ overnight
with no flow, which the raw difference would read as a large inflow. Let $k_t$ be the
split factor on date $t$, equal to $1$ on ordinary days, $2$ for a two-for-one split,
and $\tfrac12$ for a one-for-two reverse split. The mechanical change is then removed,
\begin{equation}
\label{eq:flow_split}
F_t = (S_t - k_t\,S_{t-1})\,\mathrm{NAV}_t,
\end{equation}
so a pure split contributes zero flow.

\subsection{The reported route}

A daily shares-outstanding series for ETFs is not available from free sources. Form
N-PORT supplies the flow directly instead. Every US fund files it monthly, reporting for
each of the three months in the period the total dollar value of share sales, redemptions,
and dividend reinvestment (Item B.8), the monthly total return (Item B.7), and net assets
at the period end (Item B.1). The net flow is
\begin{equation}
\label{eq:flow_reported}
F_t = \text{sales}_t - \text{redemptions}_t,
\end{equation}
with reinvestment excluded as internal recycling of distributions. This is the realized
data source: monthly granularity, with filings posted about two months after the period.
Equation \eqref{eq:flow_split} remains the route for daily data when a shares feed is
available.

\section{Normalizing}
\label{sec:normalizing}

A flow of ten million dollars is large for a small fund and negligible for a large one.
Dividing by starting size gives the organic growth rate,
\begin{equation}
\label{eq:normalized_flow}
g_t = \frac{F_t}{A_{t-1}},
\end{equation}
read as the fraction of the fund's size that arrived as new money in period $t$. Dollar
flow answers how much capital moved; $g_t$ makes funds of different sizes comparable.

Equation \eqref{eq:normalized_flow} needs $A_{t-1}$ every month, but N-PORT reports net
assets only at quarter-ends. We fill the gaps with the identity \eqref{eq:aum_identity},
\begin{equation}
\label{eq:aum_recursion}
A_t =
\begin{cases}
\widehat{A}_t & \text{if net assets } \widehat{A}_t \text{ are reported in month } t,\\[4pt]
A_{t-1}\,(1 + r_t) + F_t & \text{otherwise,}
\end{cases}
\end{equation}
snapping to each reported value and rolling \eqref{eq:aum_identity} across the two
intervening months. Snapping each quarter is required because $r_t$ is a total return:
it includes distributions, which leave the fund as cash, so an unbroken roll overstates
assets in proportion to dividend yield. Re-anchoring at every report bounds that bias to
a single quarter. The remaining sub-quarter error is small enough to leave $g_t$ accurate
(Section~\ref{sec:validation}).

\section{Aggregating to categories}
\label{sec:aggregation}

A single fund is not a sector. We assign each fund $i$ to a category $c$ and sum within it,
\begin{align}
F_{c,t} &= \sum_{i \in c} F_{i,t}, \label{eq:category_flow}\\
A_{c,t} &= \sum_{i \in c} A_{i,t}, \qquad
g_{c,t} = \frac{F_{c,t}}{A_{c,t-1}}. \label{eq:category_g}
\end{align}
The category growth rate \eqref{eq:category_g} is the asset-weighted average of its member
funds' growth rates, so larger funds count more. The market as a whole is the category
containing every fund,
\begin{equation}
\label{eq:universe_g}
g_{U,t} = \frac{\sum_i F_{i,t}}{\sum_i A_{i,t-1}},
\end{equation}
which serves as the baseline in Section~\ref{sec:rotation}.

\section{Rotation}
\label{sec:rotation}

Rotation is movement of capital between categories over time. Five statistics describe it.

Monthly flow is noisy. Trailing cumulative flow over a window of $W$ months shows the
sustained trend,
\begin{equation}
\label{eq:cumulative_flow}
\mathrm{CF}_{c,t}(W) = \sum_{k=0}^{W-1} F_{c,t-k}.
\end{equation}

Categories differ in their typical flow volatility, so raw $g_{c,t}$ is not comparable
across them. Standardizing against each category's own recent history gives a comparable
score,
\begin{equation}
\label{eq:zscore}
z_{c,t} = \frac{g_{c,t} - \mu_{L}(g_c)}{\sigma_{L}(g_c)},
\end{equation}
where $\mu_L$ and $\sigma_L$ are the mean and sample standard deviation of $g_c$ over the
$L$ months \emph{preceding} $t$. The current month is excluded from its own baseline, so a
single extreme month neither inflates $\sigma_L$ nor drags $\mu_L$ toward itself; the score
measures how surprising month $t$ is given the history before it, with no upper bound. A
value $z_{c,t}=2$ means the category is drawing unusually strong inflows relative to its
own norm.

Rotation is inherently relative: capital leaving one category to enter another. Comparing
each category to the market separates that movement from a market-wide tide,
\begin{equation}
\label{eq:relative_flow}
\mathrm{rel}_{c,t} = g_{c,t} - g_{U,t}.
\end{equation}
A positive value means the category attracts more than its market-weighted share of flow.
Asset-weighting \eqref{eq:relative_flow} and summing over categories gives zero, since the
weighted average of $g_{c,t}$ is $g_{U,t}$ by \eqref{eq:universe_g}; relative flows are a
zero-sum reallocation.

The reallocation has a measurable size. The dollar rotation flow of a category,
\begin{equation}
\label{eq:rotation_dollars}
R_{c,t} = A_{c,t-1}\,\mathrm{rel}_{c,t} = F_{c,t} - A_{c,t-1}\,g_{U,t},
\end{equation}
is its flow in excess of its assets growing at the universe rate, splitting the flow as
$F_{c,t} = A_{c,t-1}\,g_{U,t} + R_{c,t}$: tide plus rotation. The $R_{c,t}$ sum to zero
each month, so the positive legs match the negative legs dollar for dollar, and either
side is the gross rotation turnover,
\begin{equation}
\label{eq:rotation_turnover}
T_t = \frac{1}{2}\sum_c \lvert R_{c,t}\rvert, \qquad
\tau_t = \frac{T_t}{\sum_c A_{c,t-1}},
\end{equation}
the smallest total of dollars that must move between categories to turn the pro-rata
allocation $\{A_{c,t-1}\,g_{U,t}\}$ into the observed one. The rate $\tau_t$ expresses
the same quantity as a fraction of prior universe assets, comparable across a sample in
which assets grow severalfold. Both are computed on raw monthly relative flow, where the
cancellation is exact, and only in months where every category carries prior assets,
since a partial sum understates the turnover. Turnover is a lower bound on cross-sector
movement: reallocation that nets out across investors within a month is invisible. Over
the panel the median month rotates \$3.0 billion, $1.2$ percent of universe assets; the
largest readings are January 2021 (\$6.7 billion) and April 2020 ($\tau_t = 4.3$
percent).

Whether a category's relative strength is rising or fading is its momentum over $D$ periods,
\begin{equation}
\label{eq:momentum}
m_{c,t} = \mathrm{rel}_{c,t} - \mathrm{rel}_{c,t-D}.
\end{equation}
Categories differ in the volatility of their relative flow, so plotting
$\mathrm{rel}_{c,t}$ directly would let a structurally volatile sector dominate the chart
while a calm one collapses toward the origin, even when the calm sector's move is the more
unusual relative to its own norm. Monthly relative flow is also noisy, which clutters the
path a category traces. We therefore smooth relative flow with a trailing $w$-month mean,
\begin{equation}
\label{eq:rel_smoothed}
\bar{\mathrm{rel}}_{c,t}(w) = \frac{1}{w}\sum_{k=0}^{w-1}\mathrm{rel}_{c,t-k},
\end{equation}
and standardize the smoothed signal against its own recent history, as in \eqref{eq:zscore},
\begin{equation}
\label{eq:rs_ratio}
\mathrm{RS}_{c,t} =
\frac{\bar{\mathrm{rel}}_{c,t} - \mu_L(\bar{\mathrm{rel}}_c)}{\sigma_L(\bar{\mathrm{rel}}_c)},
\end{equation}
with $\mu_L$ and $\sigma_L$ taken over the $L$ months preceding $t$, and take its momentum
over $D$ periods in the sense of \eqref{eq:momentum},
\begin{equation}
\label{eq:rs_momentum}
\mathrm{RS}^{\mathrm{m}}_{c,t} = \mathrm{RS}_{c,t} - \mathrm{RS}_{c,t-D}.
\end{equation}
The rotation graph plots relative strength $\mathrm{RS}_{c,t}$ (horizontal) against its
momentum $\mathrm{RS}^{\mathrm{m}}_{c,t}$ (vertical), centered at the origin, placing each
category in one of four quadrants: leading (strong and rising), weakening (strong and
falling), lagging (weak and falling), and improving (weak and rising). Standardizing puts
every sector on a common scale, so the chart compares rotation rather than raw volatility;
the trailing mean keeps a single noisy month from dominating the path. The figures use
$w=3$, $L=12$, and $D=3$, a quarter for the smoothing and momentum horizons. A category
typically traces these quadrants clockwise, and the path is the rotation.

\section{Validation}
\label{sec:validation}

The flows are reported rather than estimated, but the reconstructed quantities admit three
independent checks. All run over the eleven-sector panel of Section~\ref{sec:data},
eighty-one months from July 2019 to March 2026.

First, N-PORT's reported monthly total return (Item B.7) is compared to the return
computed from public prices. Across the eleven sectors the correlation exceeds $0.99$,
with mean absolute error between $0.17$ and $0.80$ percentage points. This confirms the
return series and the month alignment.

Second, the asset recursion \eqref{eq:aum_recursion} is tested at the quarterly horizon.
Starting from each reported net-assets anchor, we roll flows and returns forward to the
next anchor and compare the prediction to the reported value. The median absolute error is
between $0.2$ and $1.0$ percent for every sector, confirming that flows, returns, and
reported assets are mutually consistent.

Third, relative flow satisfies a conservation identity. Weighting
\eqref{eq:relative_flow} by prior assets and summing over categories, the definition of
the baseline \eqref{eq:universe_g} gives
\begin{equation}
\label{eq:flow_conservation}
\sum_c A_{c,t-1}\,\mathrm{rel}_{c,t}
  \;=\; \sum_c F_{c,t} \;-\; g_{U,t} \sum_c A_{c,t-1} \;=\; 0,
\end{equation}
so every dollar of above-market flow in one category is matched by below-market flow
elsewhere: rotation is zero-sum in dollars by construction. The identity holds whenever
every category carries prior assets. At a category's first month the growth rate is
undefined, the category leaves the sum, and the residual equals its own flow; the
eleven-sector panel has no mid-history entries. The built panel satisfies the identity
in every month to a relative magnitude of $10^{-16}$, machine precision; a violation
indicates an aggregation or alignment defect in the panel, not a property of the data.
The raw dollar flow $\sum_c F_{c,t}$ does not net to zero and is not meant to: the
median month's net flow across the eleven sectors is $42$ percent of gross flow, the
common tide into or out of the sector complex that the relative measure removes.

A single month of the panel illustrates the output. Table~\ref{tab:snapshot} reports net
flow and growth rate by sector for March 2026.

\begin{table}[h]
\centering
\begin{tabular}{llrr}
\toprule
Ticker & Sector & Net flow (\$M) & $g$ (\%) \\
\midrule
XLE  & Energy                 & $+2{,}024$ & $+5.38$ \\
XLU  & Utilities              & $+1{,}044$ & $+4.28$ \\
XLY  & Consumer Discretionary & $+267$     & $+1.18$ \\
XLK  & Technology             & $+199$     & $+0.23$ \\
XLF  & Financials             & $+80$      & $+0.16$ \\
XLV  & Health Care            & $-912$     & $-2.11$ \\
XLC  & Communication Services & $-1{,}484$ & $-5.46$ \\
\bottomrule
\end{tabular}
\caption{Reported net flow and growth rate by sector, March 2026 (selected rows).
Capital enters Energy and Utilities and leaves Communication Services and Health Care.}
\label{tab:snapshot}
\end{table}

\section{Data and scope}
\label{sec:data}

The source is SEC Form N-PORT: monthly, public, with a reporting lag near two months, and
coverage from late 2019. The universe is the eleven Select Sector SPDR funds, one per GICS
sector, which together tile the US equity market. Each sector currently holds one fund, so
\eqref{eq:category_flow} reduces to an identity and \eqref{eq:category_g} equals that fund's
growth rate; the aggregation becomes load-bearing once a category holds several funds.

Three data effects are handled explicitly. Splits enter through the factor $k_t$ in
\eqref{eq:flow_split}. Distributions enter through the re-anchoring in
\eqref{eq:aum_recursion}, which prevents total return from inflating reconstructed assets.
The reporting lag means the most recent month is provisional until its filing posts.

ETFs are a large but partial slice of invested capital; mutual funds, separately managed
accounts, and direct holdings are outside this universe. The panel is a proxy for sector
positioning, not a complete account of it.

\appendix
\section{Notation}
\label{app:notation}

\begin{tabular}{ll}
\toprule
Symbol & Meaning \\
\midrule
$t$ & time period (month) \\
$i$ & a single fund \\
$c$ & a category (sector) \\
$U$ & the whole-market aggregate \\
$S_t$ & shares outstanding at end of $t$ \\
$\mathrm{NAV}_t$ & net asset value per share at end of $t$ \\
$A_t$ & assets under management, $S_t\,\mathrm{NAV}_t$ \\
$r_t$ & total return over $t$ \\
$k_t$ & split factor ($1$ ordinarily, $2$ for two-for-one) \\
$F_t$ & net dollar flow over $t$ \\
$g_t$ & normalized flow, $F_t / A_{t-1}$ \\
$\mathrm{CF}_{c,t}(W)$ & trailing cumulative flow over $W$ months \\
$z_{c,t}$ & standardized flow over lookback $L$ \\
$\mathrm{rel}_{c,t}$ & relative flow, $g_{c,t} - g_{U,t}$ \\
$m_{c,t}$ & flow momentum over $D$ periods \\
$\bar{\mathrm{rel}}_{c,t}(w)$ & trailing $w$-month mean of relative flow \\
$\mathrm{RS}_{c,t}$ & standardized relative flow (rotation-graph $x$) \\
$\mathrm{RS}^{\mathrm{m}}_{c,t}$ & momentum of $\mathrm{RS}$ (rotation-graph $y$) \\
\bottomrule
\end{tabular}

\end{document}
